\documentclass[11pt]{article}
\usepackage[margin=0.85in]{geometry}
\usepackage{amsmath,booktabs,graphicx,hyperref}
\setlength{\parindent}{0pt}
\setlength{\parskip}{4pt}

\title{CS 498: HW1 Report}
\author{Your Name}
\date{Fall 2026}

\begin{document}
\maketitle

\section{Rigid fitting and RANSAC}
Report SVD and RANSAC translation/rotation errors, inlier ratio, inlier RMSE,
and both alignment figures. Briefly explain the RANSAC threshold tradeoff.

\section{Real-data ICP and odometry}
Include the Open3D demo before/after alignment, ICP convergence, GT/ICP 2D
trajectory, and GT/ICP accumulated BEV. State the ICP iteration-video filename.
Also state the online odometry-video filename and explain what its current-sweep
and accumulated-intensity panels reveal about local alignment and drift. Note
how the inferred trajectory overlay grows through the current frame.
Report demo final RMSE, planar XY odometry ATE (first pose identity, no further
alignment), mean one-step translation RTE, and mean
one-step rotation RTE. Identify where drift appears in both the path and mapped
structures.

\section{Semantic voxel mapping}
First estimate the shape, voxel count, and float32 score memory for the dense
grid posed in the handout, and explain why the map should be sparse.
Include the Open3D shoulder-view cube map and state both MP4 filenames: the
turntable and the side-by-side moving LiDAR/map view. Report frames aggregated,
occupied voxels, voxel accuracy, per-class IoU for present classes, mean IoU,
observation-count statistics, and mean normalized entropy of mean log evidence.
Confirm that \texttt{voxel\_coordinates\_correct} is true. Discuss one region
with competing semantic evidence and one limitation of weighting all
observations equally. Briefly relate voxel size to detail and memory.

\section{Failure analysis}
Connect one visual failure to a specific registration or fusion assumption.

\end{document}
